\section{Metodologia}

\subsection{Arquitetura do Sistema}

\begin{figure}[H]
    \centering
    \includegraphics[width=0.49\textwidth]{figuras/diagrama_jetpack_compose_kotlin.png}
    \caption{Diagrama de blocos do projeto}
    \label{fig:diagrama_de_blocos}
\end{figure}

\subsection{Aplicativo Mobile}

O aplicativo mobile foi desenvolvido com Kotlin Multiplatform e Compose Multiplatform, adotando o módulo compartilhado \texttt{composeApp} para a interface, modelos de domínio, componentes visuais, contratos de sessão remota e parte da lógica de jogo. As integrações dependentes de plataforma permanecem isoladas em implementações específicas. Na versão atual do projeto, o Android é a plataforma funcional de referência, concentrando autenticação, navegação autenticada, coleta real de sensores, permissões, execução em primeiro plano, comunicação HTTP, conexão WebSocket/STOMP e integração com Wear OS. O iOS permanece como um invólucro multiplataforma por meio do \texttt{MainViewController()} hospedado pelo aplicativo SwiftUI, mas ainda sem paridade com o fluxo Android de telemetria, sessão online e renderização KorGE.

A aplicação funciona como o ponto de contato do usuário com a sessão gamificada. Ela realiza autenticação e cadastro, carrega as Hordas disponibilizadas pelo backend, permite selecionar a Horda da corrida, coleta sinais de movimento e biofeedback, acompanha o estado remoto da sessão, renderiza localmente a cena do jogo e envia dados biométricos ao backend em tempo real. A raiz compartilhada \texttt{App.kt} aplica o tema visual, enquanto a orquestração completa do fluxo Android ocorre em \texttt{MainActivity}, que conecta autenticação, navegação, telemetria, sessão remota e ponte com smartwatch.

A interface do aplicativo é construída com Compose Multiplatform e organizada por funcionalidades. As telas de autenticação permitem login e cadastro de usuários; a tela inicial apresenta as Hordas carregadas do backend, a seleção da Horda e o início da corrida; a tela de jogo exibe a cena gamificada, o progresso remoto, o estado da sessão, os sinais de telemetria e os indicadores de desempenho; a tela de ranking apresenta a classificação da sessão; e as telas de perfil permitem visualizar e editar dados do usuário. Também existem telas auxiliares para permissões de telemetria, histórico e estados de conexão com o relógio.

\vspace{0.2cm}
\subsubsection{Pipeline de telemetria}
\vspace{0.2cm}

O pipeline de telemetria é uma das partes centrais do aplicativo. No Android, a entrada dos sinais ocorre por meio do \texttt{FusedLocationTrackingService}, que obtém localização via Play Services Location, do \texttt{AndroidMotionSensorService}, que lê amostras de aceleração do dispositivo, e do \texttt{AndroidWearTelemetryBridge}, que recebe mensagens do smartwatch via Wear OS. As mensagens enviadas no caminho \texttt{/telemetry} são processadas pelo \texttt{WearTelemetryMessageParser} e convertidas em amostras de biofeedback com frequência cardíaca, permitindo que o estado de telemetria incorpore dados reais do wearable. A aplicação também possui o \texttt{WearTelemetryListenerService}, que recebe mensagens do relógio mesmo fora do listener registrado diretamente pelo bridge.

Essas fontes são centralizadas pelo \texttt{DefaultTelemetryRepository}, que expõe um \texttt{StateFlow<TelemetryState>} para o restante da aplicação. O repositório controla o ciclo de vida da sessão de movimento, monitora disponibilidade de sensores e permissões, agrega localização, aceleração e possíveis sinais biométricos, calcula a distância acumulada a partir dos pontos de localização e publica o estado atualizado da telemetria. A disponibilidade considera permissão de localização, provedor de localização ativo, sensor de movimento e presença do relógio, registrando problemas no próprio \texttt{TelemetryState}.

O estado publicado mantém a sessão de movimento, a última localização, a última aceleração, a última amostra de biofeedback e a estratégia de telemetria escolhida. O \texttt{SelectTelemetryStrategyUseCase} alterna entre \texttt{MOVEMENT\_ONLY} e \texttt{BPM\_AND\_MOVEMENT} conforme a conexão do relógio e a disponibilidade de BPM. Na integração online atual, esses dados são convertidos pelo \texttt{buildBiometricDataMessage} em uma mensagem com BPM, velocidade, pace e distância acumulada, que é enviada ao backend durante a sessão.

\vspace{0.2cm}
\subsubsection{Lógica de jogo e gamificação}
\vspace{0.2cm}

A gamificação foi modelada como uma simulação de corrida e fuga. O usuário representa o corredor, enquanto a Horda atua como adversário virtual que pressiona o jogador ao longo da sessão. O domínio do jogo não acessa diretamente sensores Android; ele recebe estados e contratos abstraídos, preservando a separação entre coleta nativa, comunicação remota e representação visual.

No aplicativo, o \texttt{GameController} concentra o estado necessário para a interface e para a renderização local da cena, expondo um \texttt{GameSnapshot} com distância, velocidade, pressão, risco, pontuação, tempo decorrido e resultado exibidos ao usuário. Em uma sessão online, os cálculos autoritativos da corrida ficam no backend: posição do usuário, posição da Horda, distância restante, velocidade, progresso e resultado são recebidos pelo aplicativo por WebSocket e convertidos em \texttt{GameSnapshot} por \texttt{RemoteGameStateSnapshotMapper}. Assim, o mobile atua principalmente como coletor de telemetria e cliente de visualização, enquanto o servidor permanece como fonte de verdade da sessão.

A visualização do jogo é renderizada com KorGE no Android. A cena principal, definida em \texttt{MainScene}, monta o fundo com parallax, o corredor e as unidades da Horda. A cada atualização, a cena lê o \texttt{GameSnapshot} mantido pelo controlador, que pode refletir o estado remoto recebido do backend, e ajusta os elementos visuais de acordo com o estado lógico da sessão. Dessa forma, a camada visual apenas representa o estado da corrida, sem criar uma fonte paralela de verdade para a simulação.

\vspace{0.2cm}
\subsubsection{Integração em tempo real com o backend}
\vspace{0.2cm}

A comunicação com o backend combina chamadas REST e WebSocket com STOMP. As chamadas REST são usadas para autenticação, cadastro, carregamento de Hordas, carregamento e edição do perfil, criação de sessão, consulta do leaderboard e encerramento da sessão. Após a autenticação, o token JWT retornado pelo backend é persistido localmente pelo gerenciador de sessão Android e enviado nas requisições protegidas.

Ao iniciar uma corrida, o aplicativo exige uma Horda selecionada e solicita ao backend a criação de uma sessão por meio do endpoint \texttt{/sessions/start}, enviando o identificador da Horda no corpo da requisição. Em seguida, o \texttt{OnlineSessionRepository} usa o \texttt{StompWebSocketClient} para abrir conexão no endpoint \texttt{/ws}, assinar os tópicos \texttt{/topic/session/\{sessionId\}/leaderboard} e \texttt{/topic/session/\{sessionId\}/game-state}, e receber atualizações de ranking e estado remoto do jogo.

Durante a sessão, o \texttt{OnlineSessionRepository} lê o \texttt{TelemetryState} local e envia mensagens biométricas para o destino STOMP \texttt{/app/train/data}. O envio é limitado a uma atualização por segundo, reduzindo tráfego desnecessário e mantendo frequência suficiente para atualização do ranking e do estado remoto em tempo real. Ao finalizar a corrida, o aplicativo encerra a coleta local, desconecta o WebSocket e solicita ao backend o fechamento da sessão pelo endpoint \texttt{/sessions/\{sessionId\}/finish}. Em paralelo, o \texttt{WatchGameProgressBridge} publica o progresso do jogo para o smartwatch no caminho \texttt{/game-progress}, incluindo distância, meta, progresso, BPM, risco, pressão da Horda e resultado local.

\vspace{0.2cm}
\subsubsection{Permissões e execução em primeiro plano}
\vspace{0.2cm}

Como a sessão depende de sinais físicos do dispositivo, a versão Android verifica permissões, conectividade e disponibilidade antes de iniciar a coleta. O aplicativo exige permissão de localização, autorização de notificação em versões recentes do Android, provedor de localização ativo, sensor de movimento disponível e conexão validada com a internet para iniciar a sessão online. Caso algum requisito de telemetria esteja ausente, o usuário é direcionado para a tela de permissões de telemetria; se a conexão cair durante uma corrida, o jogo é encerrado localmente e a sessão remota é marcada com erro.

Durante a corrida, o \texttt{TelemetryForegroundService} mantém a sessão ativa em primeiro plano e atualiza a notificação associada ao acompanhamento. Essa estratégia evita que a coleta seja interrompida de forma prematura pelo sistema operacional durante uma sessão em andamento, ao mesmo tempo em que mantém transparência para o usuário sobre a execução da telemetria.

\vspace{0.2cm}
\subsubsection{Qualidade e testes no mobile}
\vspace{0.2cm}

A base de testes está concentrada principalmente no código compartilhado e em componentes Android específicos. Há testes para o mapeamento de estado remoto para \texttt{GameSnapshot}, para a configuração de Hordas como \texttt{SessionConfig}, para o estado de sessão remota, para o mapeamento de telemetria em mensagens biométricas, para o mapeamento visual de distância, para modelos locais de Horda, para formatação de notificação de telemetria e para formatação de campos de interface. Também existem testes unitários Android para o cliente STOMP, decodificação de mensagens de leaderboard, montagem de requisições HTTP e testes instrumentados voltados a componentes de interface.

Essa distribuição de testes prioriza os pontos de maior risco arquitetural da versão atual: compatibilidade entre o aplicativo mobile e os contratos expostos pelo backend, conversão de estados remotos em representação local do jogo, formatação de dados enviados pela sessão online e estabilidade de componentes reutilizáveis de interface.

\subsection{Backend}

O \textit{backend} da aplicação foi desenvolvido com Spring Boot 3.4.4 e Kotlin 2.1.20, seguindo uma arquitetura em camadas que separa responsabilidades entre apresentação, lógica de negócio e persistência. A escolha por Kotlin sobre a JVM permite aproveitar a maturidade do ecossistema Spring ao mesmo tempo que oferece concisão sintática e \textit{null safety} em nível de linguagem, características relevantes para a robustez de uma API com múltiplos fluxos de dados concorrentes.

\vspace{0.2cm}
\noindent\textbf{Organização em camadas.}
\vspace{0.2cm}

A estrutura interna do \textit{backend} é organizada em quatro camadas horizontais, cada uma com responsabilidade bem delimitada.

A \textbf{camada de configuração} reúne componentes transversais à aplicação: o filtro de autenticação JWT, a política de segurança HTTP, a configuração do \textit{broker} de mensagens WebSocket e o tratamento centralizado de exceções. Por concentrar essas preocupações em um único pacote (\texttt{config}), a camada evita que lógica de segurança e de infraestrutura se espalhe pelos demais componentes.

A \textbf{camada de apresentação} expõe os \textit{endpoints} REST sob os prefixos \texttt{/auth}, \texttt{/sessions} e \texttt{/users}, além do \textit{handler} WebSocket responsável por receber o fluxo contínuo de dados biométricos. Toda a documentação interativa dos \textit{endpoints} é gerada automaticamente via SpringDoc OpenAPI, acessível pela interface Swagger integrada à aplicação.

A \textbf{camada de serviço} concentra toda a lógica de negócio: autenticação e geração de tokens JWT, ciclo de vida das sessões de treino, cálculo da posição virtual do competidor automático denominado Horda, e gerenciamento do estado da sessão em \textit{cache}. Os serviços de acesso ao Redis foram deliberadamente separados dos demais serviços em um subpacote próprio (\texttt{service/redis}), evidenciando a distinção entre operações de persistência durável e operações de \textit{cache} em memória.

A \textbf{camada de domínio} contém as entidades JPA mapeadas para o banco de dados relacional (\texttt{User}, \texttt{TrainSession}, \texttt{BiometricData}, \texttt{Horde}, \texttt{Ranking}, \texttt{Achievement} e \texttt{Friendship}), cada uma acompanhada de seu respectivo repositório Spring Data. O esquema relacional é versionado e aplicado de forma incremental pelo Flyway, garantindo rastreabilidade das migrações ao longo do desenvolvimento.

\vspace{0.2cm}
\noindent\textbf{Estratégia de persistência híbrida.}
\vspace{0.2cm}

Uma decisão arquitetural central deste trabalho é o uso combinado de dois mecanismos de persistência com papéis distintos: Redis para o estado volátil de alta frequência, e PostgreSQL para a persistência durável ao final de cada sessão.

Durante a execução de uma sessão de treino, nenhuma operação de escrita é realizada no banco relacional. Cada atualização biométrica recebida pelo servidor aciona exclusivamente operações Redis, que completam em menos de 1 ms por chamada. O \textit{leaderboard} é mantido como um \textit{Sorted Set} (ZSET), estrutura de dados nativa do Redis que indexa os participantes pela distância acumulada e permite atualizar a posição de um usuário em tempo $O(\log N)$ via operação \texttt{ZADD}, bem como recuperar os primeiros colocados em $O(\log N + N)$ via \texttt{ZREVRANGE}. Paralelamente, o estado biométrico instantâneo de cada participante (frequência cardíaca, cadência, pace, calorias) é armazenado em um \textit{Hash} Redis por usuário, permitindo leituras e atualizações de campos individuais sem reescrever o objeto completo.

Ao encerrar a sessão, o \texttt{TrainSessionService} lê o \textit{leaderboard} final do Redis, persiste os rankings individuais no PostgreSQL como registros da entidade \texttt{Ranking}, registra a distância total percorrida na entidade \texttt{TrainSession} e, por fim, reduz o TTL das chaves Redis para uma hora. Essa janela mantém os dados disponíveis para consultas imediatas pós-corrida sem ocupar memória de forma indefinida.

Essa separação garante que o caminho crítico do processamento de cada mensagem biométrica recebida em tempo real nunca seja bloqueado por operações de I/O no banco relacional, mantendo a latência de ponta a ponta abaixo de 10 ms por mensagem.

\vspace{0.2cm}
\noindent\textbf{Segurança e autenticação.}
\vspace{0.2cm}

A aplicação adota autenticação \textit{stateless} baseada em JWT (\textit{JSON Web Token}). Ao realizar o login, o usuário recebe um token assinado com algoritmo HMAC-SHA, que deve ser enviado no cabeçalho \texttt{Authorization} de todas as requisições subsequentes. O \texttt{JwtAuthFilter}, implementado como \texttt{OncePerRequestFilter}, intercepta cada requisição HTTP, extrai e valida o token antes que ela chegue aos \textit{controllers}. As senhas dos usuários são armazenadas com \textit{hash} BCrypt, eliminando a exposição de credenciais em texto plano.

O servidor não mantém estado de sessão HTTP, o que simplifica o escalonamento horizontal da aplicação, já que qualquer instância pode validar um token de forma independente sem necessidade de sessões compartilhadas.

\vspace{0.2cm}
\noindent\textbf{Infraestrutura e qualidade.}
\vspace{0.2cm}

Todo o ambiente da aplicação é orquestrado via Docker Compose, simplificando a reprodução do ambiente de desenvolvimento e de testes. Um \textit{pipeline} de integração contínua configurado no GitHub Actions executa automaticamente a suíte de testes a cada \textit{push}, gerando relatórios de cobertura por meio do \textit{plugin} JaCoCo. Os testes cobrem todas as camadas da aplicação utilizando o banco H2 em memória para isolamento dos testes de persistência relacional.

\subsection{Protocolos de Comunicação}

O protocolo WebSocket, combinado com o STOMP (\textit{Simple Text Oriented Messaging Protocol}) como camada de mensageria, viabiliza a comunicação bidirecional e assíncrona entre o aplicativo móvel e o servidor. Ao contrário do modelo HTTP tradicional, em que o cliente deve realizar requisições periódicas para verificar atualizações (\textit{polling}), o WebSocket mantém uma conexão persistente que permite ao servidor enviar dados ao cliente de forma proativa.

O aplicativo mobile, desenvolvido em Kotlin Multiplatform com interface em Compose, conecta-se ao \textit{endpoint} \texttt{/ws} no momento em que o usuário inicia a sessão. A partir desse ponto, os dados biométricos e de movimento processados localmente são encaminhados ao destino STOMP \texttt{/app/train/data}, onde o backend os recebe, processa e persiste no Redis. Após cada atualização, o servidor transmite o \textit{leaderboard} recalculado para todos os clientes inscritos no tópico \texttt{/topic/session/\{sessionId\}/leaderboard}, possibilitando que todos os participantes visualizem o ranking atualizado simultaneamente e em tempo real.

\subsection{Lógica da Gamificação}

A Horda representa um adversário virtual que avança de forma autônoma durante a sessão de treino. Sua posição é calculada pelo \texttt{HordePositionService} com base no tempo decorrido desde o início da sessão e no \textit{pace}-alvo configurado, em minutos por quilômetro. Como o tempo de início da sessão é lido diretamente do Redis, o cálculo não exige consulta adicional ao banco relacional e permanece puramente aritmético. O modelo de domínio prevê ainda sub-hordas que herdam o \textit{pace} da Horda pai, permitindo criar variações de dificuldade sem duplicar configurações.
